% =====================================================
%  DAY 2 -- PRACTICE CODE SHEET
%  Math Mode for Economists
%  "From \begin to \end: The Complete LaTeX Guide"
%  White Globe Initiative | Instructor: Atikur Rahman
%
%  Every block below is exactly the code taught in the Day 2
%  slides, assembled into one file so you can compile every
%  example's output at once. Compile with pdflatex.
% =====================================================

\documentclass[11pt,a4paper]{article}
\usepackage[utf8]{inputenc}
\usepackage{amsmath,amssymb}  % Topic: The amsmath Package

\begin{document}

% ---- Topic: Inline vs. Display Math ----
\section*{Inline vs.\ Display Math}
The elasticity is $\varepsilon = \frac{dQ}{dP}\cdot\frac{P}{Q}$.

\begin{equation}
  \pi = TR - TC
\end{equation}

% ---- Topic: Subscripts and Superscripts ----
\section*{Subscripts and Superscripts}
$x_i$              % single-character subscript
\quad
$Y_{it}$           % braces needed for multi-character
\quad
$X_{i,t-1}^{2}$    % subscript AND superscript, lagged

% ---- Topic: Greek Letters and Common Symbols ----
\section*{Greek Letters and Common Symbols}
$\alpha$  $\beta$  $\theta$  $\lambda$  $\sigma$
\quad
$\Sigma$  $\Delta$  $\Omega$  $\infty$

% ---- Topic: Operators: Sums, Products, Limits, Max ----
\section*{Operators: Sums, Products, Limits, Max}
$\displaystyle\sum_{i=1}^{n} x_i \qquad \lim_{x\to\infty} \qquad
\operatorname*{argmax}_{x} f(x)$

% ---- Topic: Fractions and Binomials ----
\section*{Fractions and Binomials}
$\varepsilon_d = \dfrac{\%\,\Delta Q_d}{\%\,\Delta P}
\qquad \dbinom{n}{k} = \dfrac{n!}{k!(n-k)!}$

% ---- Topic: Multi-line Derivations with align ----
\section*{Multi-line Derivations with align}
\begin{align}
  \hat{\beta} &= (X'X)^{-1}X'Y \\
  \text{Var}(\hat{\beta}) &= \sigma^2 (X'X)^{-1}
\end{align}

% ---- Topic: Matrices: Linear Algebra & Game Theory ----
\section*{Matrices: Linear Algebra \& Game Theory}
$\begin{pmatrix} 1 & 0.5 \\ 0.5 & 1 \end{pmatrix}
\qquad
\begin{pmatrix} (3,3) & (0,5) \\ (5,0) & (1,1) \end{pmatrix}$

% ---- Topic: Labeling and Referencing Equations ----
\section*{Labeling and Referencing Equations}
\begin{equation}
  \label{eq:cobb-douglas}
  Y = A K^{\alpha} L^{1-\alpha}
\end{equation}

As shown in Equation~\eqref{eq:cobb-douglas}, output depends on...

\end{document}
